\documentclass[11pt,a4paper]{article}
\usepackage[T1]{fontenc}
\usepackage{lmodern,microtype}
\usepackage[margin=25mm]{geometry}
\usepackage{amsmath,amssymb,amsthm,booktabs,array}
\usepackage[hidelinks]{hyperref}
\ifdefined\pdfinfoomitdate\pdfinfoomitdate=1\fi
\ifdefined\pdftrailerid\pdftrailerid{}\fi
\ifdefined\pdfsuppressptexinfo\pdfsuppressptexinfo=15\fi
\hypersetup{pdftitle={Sun's supercongruence for an Apery-like sequence},
 pdfauthor={Oleksiy Babanskyy},
 pdfsubject={A binomial supercongruence, modular periods and a finite differential invariant},
 pdfkeywords={supercongruences, Apery-like sequences, Franel numbers, modular forms, complex multiplication}}
\pagestyle{plain}
\newtheorem{theorem}{Theorem}[section]
\newtheorem{proposition}[theorem]{Proposition}
\newtheorem{lemma}[theorem]{Lemma}
\newtheorem{corollary}[theorem]{Corollary}
\theoremstyle{remark}\newtheorem{remark}[theorem]{Remark}
\newcommand{\Z}{\mathbb Z}\newcommand{\Q}{\mathbb Q}
\newcommand{\C}{\mathbb C}\newcommand{\Fp}{\mathbb F_p}
\newcommand{\HH}{\mathfrak H}
\newcommand{\A}{\mathcal A}\newcommand{\B}{\mathcal B}
\newcommand{\dd}{\partial}\newcommand{\Tp}{\mathcal T_p}
\newcommand{\eps}{\epsilon_p}
\newcommand{\leg}[2]{\left(\frac{#1}{#2}\right)}
\newcommand{\smat}[4]{\begin{pmatrix}#1&#2\\#3&#4\end{pmatrix}}
\DeclareMathOperator{\ord}{ord}
\allowdisplaybreaks
\title{Sun's supercongruence for an Ap\'ery-like sequence}
\author{Oleksiy Babanskyy}
\date{}
\begin{document}
\maketitle
\begin{abstract}
Let $a_n=\sum_{j=0}^n\binom nj^2\binom{2j}j$.
We prove that, for every prime $p>3$, the sum
\[
 (-1)^{(p-1)/2}\sum_{n=0}^{p-1}\binom{2n}n a_n/54^n
\]
is congruent modulo $p^2$ to $4x^2-2p$ when
$p=x^2+3y^2\equiv1\pmod3$, and to zero when $p\equiv2\pmod3$.
For the latter primes a quadratic differential invariant of the truncated
generating polynomial lifts Sun's congruence modulo $p$.
For the former primes we certify two modular periods directly by identities
of weight eight, construct a Hecke polynomial with an explicit degree bound,
and specialize a finite polynomial identity at a complex multiplication point.
The coefficient at degree $p$ and both parity cases of the complex
multiplication matrices are treated explicitly. A corresponding evaluation
of the Franel sum with denominator $50$ follows by the same arguments.
\end{abstract}

\section{The supercongruence and its two branches}\label{sec:intro}
For $n\ge0$ define
\begin{equation}\label{eq:sequences}
 a_n=\sum_{j=0}^n\binom nj^2\binom{2j}j,
 \qquad f_n=\sum_{j=0}^n\binom nj^3,
 \qquad g_n=\binom{2n}n a_n,
 \qquad b_n=\binom{2n}n f_n.
\end{equation}
The integers $f_n$ are the Franel numbers. Write
\begin{equation}\label{eq:sums}
 G(t)=\sum_{n\ge0}g_nt^n,\qquad F(t)=\sum_{n\ge0}b_nt^n,
 \qquad Y_p(t)=\sum_{n=0}^{p-1}y_nt^n
 \quad\text{if }Y(t)=\sum_{n\ge0}y_nt^n.
\end{equation}
In particular, all truncations in this article end at $p-1$. We use
\begin{equation}\label{eq:B-T}
 B_p=G_p(1/54),\qquad T_{50,p}=F_p(1/50),
 \qquad\eps=(-1)^{(p-1)/2}.
\end{equation}
The first sum is $p$-integral for $p>3$; the second is used only for $p>5$.
Congruences between $p$-integral rational numbers are interpreted in $\Z_p$.

\begin{theorem}\label{thm:main}
For every prime $p>3$,
\begin{equation}\label{eq:main}
 \eps B_p\equiv
 \begin{cases}
 4x^2-2p,&p=x^2+3y^2\equiv1\pmod3,\\
 0,&p\equiv2\pmod3
 \end{cases}
 \pmod{p^2}.
\end{equation}
\end{theorem}

This is Conjecture~3.5 of Zhi-Hong Sun
\cite[p.~2742]{sun}. Every prime $p\equiv1\pmod3$ admits the indicated
representation; the signs of $x,y$ do not affect the expression $4x^2-2p$.
Theorem~3.5 of the same paper, on p.~2739, concerns the distinct congruence
modulo $p$. Its vanishing statement and proof give the input
\begin{equation}\label{eq:Sun-input}
 B_p\equiv T_{50,p}\equiv0\pmod p
 \qquad(p>5,\ p\equiv2\pmod3).
\end{equation}
Only this modulo-$p$ result is used in the inert branch below.

The proof separates primes inert and split in $\Q(\sqrt{-3})$.
Section~\ref{sec:invariant} proves a finite differential invariant and
uses it to lift \eqref{eq:Sun-input}. Sections~\ref{sec:periods}
and~\ref{sec:geometry} identify modular periods, their divisors and an exact
complex multiplication value. Section~\ref{sec:Hecke} constructs a Hecke
polynomial and proves a polynomial identity modulo $p^2$.
Section~\ref{sec:split} computes a unit root in both split parity cases and
completes the proof. The modular polynomial construction follows the method
of Beukers \cite[Lemma~3.7 and Section~5]{beukers}; all ingredients needed
here at modulus $p^2$ are proved explicitly. The directly certified periods
also identify the Franel parametrization recorded in
\cite[Section~9.15]{beukers}.

\section{A finite differential invariant and the inert primes}
\label{sec:invariant}
We first work with a general integral sequence $u_n$, parameters
$a,b,c\in\Z$, and recurrence
\begin{equation}\label{eq:recurrence}
 u_0=1,\quad u_1=b,\qquad
 (n+1)^2u_{n+1}=(an(n+1)+b)u_n-cn^2u_{n-1}\quad(n\ge1).
\end{equation}
The sequences $a_n$ and $f_n$ satisfy this recurrence with parameters
$(10,3,9)$ and $(7,2,-8)$ respectively
\cite[pp.~2729--2730]{sun}. Exact telescoping proofs are provided in
Appendix~\ref{app:recurrences}. Let
$v_n=\binom{2n}n u_n$, $V(t)=\sum_{n\ge0}v_nt^n$ and $\theta=t\,d/dt$.
The ratios of successive central binomial coefficients give
\begin{equation}\label{eq:central-rec}
 n^3v_n=2(2n-1)(an(n-1)+b)v_{n-1}
 -4c(n-1)(2n-3)(2n-1)v_{n-2}\quad(n\ge1),
\end{equation}
where $v_{-1}=0$. Hence $LV=0$, with
\begin{equation}\label{eq:L}
 L=\theta^3-2t(2\theta+1)(a\theta^2+a\theta+b)
       +4ct^2(\theta+1)(2\theta+1)(2\theta+3).
\end{equation}
Put $D(t)=1-4at+16ct^2$. With coefficients to the left of the operators,
\begin{align}\label{eq:L-expanded}
 L={}&D\theta^3+(-6at+48ct^2)\theta^2\notag\\
 &+(-2(a+2b)t+44ct^2)\theta-2bt+12ct^2.
\end{align}
Its leading coefficient in ordinary derivatives is $t^3D(t)$.

For any polynomial or formal series $Y$ define
\begin{equation}\label{eq:invariant}
 I(Y)=D\bigl(2Y\theta^2Y-(\theta Y)^2\bigr)
       +(\theta D)Y\theta Y+4t(-b+3ct)Y^2.
\end{equation}
\begin{lemma}\label{lem:differential}
The identity
\begin{equation}\label{eq:Leibniz}
 \theta I(Y)=2YLY
\end{equation}
holds for every $Y$. In particular $I(V)=0$ in $\Q[[t]]$.
\end{lemma}
\begin{proof}
Set $E=4t(-b+3ct)$. Applying the product rule to \eqref{eq:invariant}
gives
\[
 2DY\theta^3Y+3(\theta D)Y\theta^2Y
 +(\theta^2D+2E)Y\theta Y+(\theta E)Y^2.
\]
This is $2YLY$ by \eqref{eq:L-expanded}. Since $LV=0$, the series $I(V)$
is constant in characteristic zero. Its constant term is zero, which proves
the last assertion.
\end{proof}

\begin{proposition}\label{prop:finite-invariant}
Suppose all $u_n$ in \eqref{eq:recurrence} are integers. For any prime
$p>3$, the polynomial $U(t)=\sum_{n=0}^{p-1}v_nt^n$ satisfies
\begin{equation}\label{eq:finite-invariant}
 I(U)\equiv0\pmod{p^2}\qquad\text{in }\Z[t].
\end{equation}
\end{proposition}
\begin{proof}
The exact residual of the truncation, computed from \eqref{eq:central-rec}, is
\begin{equation}\label{eq:LU}
 LU=-p^3v_pt^p+4cp(2p-1)(2p+1)v_{p-1}t^{p+1}.
\end{equation}
All coefficients below degree $p$ vanish, the omitted $v_p$ supplies the
first term, and the highest possible remaining degree is $p+1$.
Since $p\mid\binom{2p-2}{p-1}$ and $u_{p-1}$ is integral,
$p\mid v_{p-1}$. Thus $LU\equiv0\pmod{p^2}$ and
$\theta I(U)\equiv0\pmod{p^2}$.

The latter conclusion alone cannot eliminate coefficients whose degrees
are divisible by $p$. Here $\deg I(U)\le2p$ and $[t^0]I(U)=0$.
All positive degrees other than $p,2p$ are units modulo $p^2$, so it remains
to calculate precisely those two coefficients over $\Z$.
Since $V-U=v_pt^p+O(t^{p+1})$ and $I(V)=0$, comparison at degree $p$
gives
\begin{equation}\label{eq:resonance-p}
 [t^p]I(U)=-2p^2v_p.
\end{equation}
Indeed the only contribution of the omitted term at this degree pairs
the constant $V(0)=1$ with $\theta^2(v_pt^p)$ in $2V\theta^2V$;
the other terms contain an extra $t$ or a derivative of the constant.
At degree $2p$, only the highest term of each factor $U$ contributes.
The resulting scalar is
\[
 16c(p-1)^2+32c(p-1)+12c=4c(2p-1)(2p+1),
\]
so
\begin{equation}\label{eq:resonance-2p}
 [t^{2p}]I(U)=4c(2p-1)(2p+1)v_{p-1}^2.
\end{equation}
Both coefficients are divisible by $p^2$. This treats every degree.
\end{proof}

\begin{theorem}[Lifting at a nonsingular parameter]\label{thm:lift}
Under the hypotheses of Proposition~\ref{prop:finite-invariant}, let
$t_0\in\Z_p$ and suppose $t_0D(t_0)\in\Z_p^\times$.
Then
\begin{equation}\label{eq:lift}
 U(t_0)\equiv0\pmod p\quad\Longrightarrow\quad
 U(t_0)\equiv0\pmod{p^2}.
\end{equation}
The conclusion holds for every $p$-integral lift of the same residue class.
\end{theorem}
\begin{proof}
Let $\overline U$ be the reduction in $\Fp[t]$. It is nonzero, since
$U(0)=1$, and has degree less than $p$. At a zero $\bar t_0$ of
$\overline U$, the reduction of \eqref{eq:finite-invariant} gives
$-D(\bar t_0)(\theta\overline U(\bar t_0))^2=0$.
Thus $\theta\overline U(\bar t_0)=0$, and $\bar t_0\ne0$ implies that
the zero has multiplicity at least two.

It has multiplicity exactly two. Otherwise write
$\overline U(t)=w(t-\bar t_0)^m+O((t-\bar t_0)^{m+1})$ with
$w\ne0$ and $3\le m\le p-1$. Equation~\eqref{eq:LU} implies
$\overline L\overline U=0$. The coefficient of local degree $m-3$ in
this equation is
\[
 \bar t_0^3D(\bar t_0)m(m-1)(m-2)w\ne0.
\]
Lower derivatives have order at least $m-2$ and cannot cancel it.
This contradiction proves that $\theta^2U(t_0)$ is a unit, since its
reduction is $\bar t_0^2\overline U''(\bar t_0)\ne0$.

Now $U(t_0)$ and $\theta U(t_0)$ are both divisible by $p$.
Evaluating the finite identity \eqref{eq:finite-invariant} at $t_0$
therefore leaves
\[
 2D(t_0)\theta^2U(t_0)U(t_0)\equiv0\pmod{p^2}.
\]
The coefficient of $U(t_0)$ is a unit, giving \eqref{eq:lift}.
All nonvanishing conditions depended only on the residue class of $t_0$.
\end{proof}

For $a_n$ and $f_n$ the discriminants specialize to
\begin{align}\label{eq:discriminants}
 D_a(t)&=(1-4t)(1-36t),&D_a(1/54)&=25/81,\notag\\
 D_f(t)&=(1+4t)(1-32t),&D_f(1/50)&=243/625.
\end{align}
The evaluation points $1/54$ and $1/50$ and the displayed discriminant
values are units for $p>5$.
Equations~\eqref{eq:Sun-input} and~\eqref{eq:lift} give
\begin{equation}\label{eq:inert-result}
 B_p\equiv T_{50,p}\equiv0\pmod{p^2}
 \qquad(p>5,\ p\equiv2\pmod3).
\end{equation}
For $p=5$ the first discriminant is not a unit, so we calculate directly.
The initial values $a_0,\ldots,a_4$ are $1,3,15,93,639$.
The five summands of $B_5$ are respectively $1,14,15,15,5$ modulo $25$;
their sum is $50\equiv0\pmod{25}$.
This completes the inert branch of Theorem~\ref{thm:main}.

\section{Direct certification of the two modular periods}
\label{sec:periods}
Let $q=e^{2\pi i\tau}$, $\tau\in\HH$, and let
$\eta_d=\eta(d\tau)$, where
$\eta(\tau)=q^{1/24}\prod_{m\ge1}(1-q^m)$.
Define
\begin{equation}\label{eq:eta-ABJ}
 \A=\frac{(\eta_2\eta_3)^7}{(\eta_1\eta_6)^5},\qquad
 \B=\frac{(\eta_1\eta_6)^7}{(\eta_2\eta_3)^5},\qquad
 J=\A\B=(\eta_1\eta_2\eta_3\eta_6)^2,
\end{equation}
and
\begin{equation}\label{eq:Theta-t}
 \Theta_+=\A+\B,\quad\Theta_-=\A-\B,
 \qquad t_+=J/\Theta_+^2,\quad t_-=J/\Theta_-^2.
\end{equation}
These products give
$\Theta_\pm\in1+q\Z[[q]]$ and $t_\pm\in q+q^2\Z[[q]]$.
The eta criterion and cusp-order formula
\cite[p.~2]{rousewebb} show that $\A,\B$ are holomorphic forms of weight
two with trivial character on $\Gamma_0(6)$, and $J$ is a form of weight four.
At the four cusps, represented by denominators $1,2,3,6$, the respective
orders are
\begin{equation}\label{eq:cusp-orders}
 \ord\A=(0,1,1,0),\qquad\ord\B=(1,0,0,1),
 \qquad\ord J=(1,1,1,1).
\end{equation}
The underlying eta exponent data are in Appendix~\ref{app:certificates}.
Eta has no zeros in $\HH$. In particular $J$ has no interior zeros, and
each $\Theta_\pm$ has a nonzero constant term at every cusp: precisely
one of $\A,\B$ has order zero there.

Distinguish $\dd=q\,d/dq=(2\pi i)^{-1}d/d\tau$ from $\theta=t\,d/dt$.
We simultaneously use the two rows
\begin{equation}\label{eq:period-rows}
 (a,b,c,\Theta)=(10,3,9,\Theta_+),\qquad
 (a,b,c,\Theta)=(7,2,-8,\Theta_-).
\end{equation}
Put $\ell=\dd J/J$ and $W=\Theta\ell-2\dd\Theta$.
\begin{proposition}\label{prop:P}
For each row of \eqref{eq:period-rows},
\begin{align}
 W^2-\Theta^4+4aJ\Theta^2-16cJ^2&=0,\label{eq:P1}\\
 2\Theta\dd^2\Theta-3(\dd\Theta)^2
       -4bJ\Theta^2+12cJ^2&=0.\label{eq:P2}
\end{align}
\end{proposition}
\begin{proof}
First we verify the space containing the two residuals.
For $\gamma=\smat{\alpha}{\beta_0}{\gamma_0}{\delta_0}\in\Gamma_0(6)$,
write $j=\gamma_0\tau+\delta_0$ and $\lambda=\gamma_0/(2\pi i)$.
Differentiating the weight-two and weight-four transformation laws yields
\begin{align*}
 \Theta(\gamma\tau)&=j^2\Theta,\\
 (\dd\Theta)(\gamma\tau)&=j^4\dd\Theta+2\lambda j^3\Theta,\\
 (\dd^2\Theta)(\gamma\tau)&=j^6\dd^2\Theta
             +6\lambda j^5\dd\Theta+6\lambda^2j^4\Theta,\\
 \ell(\gamma\tau)&=j^2\ell+4\lambda j.
\end{align*}
Consequently $W$ transforms with weight four, and
$2\Theta\dd^2\Theta-3(\dd\Theta)^2$ transforms with weight eight.
Every other term in \eqref{eq:P1} and~\eqref{eq:P2} has weight eight.
There are no interior poles, since the only denominator at this stage is $J$.

At a cusp choose a scaling matrix $\sigma$, and put
$y=\Theta|_2\sigma$, $j_0=J|_4\sigma$.
The same differentiation laws give
\[
 W|_4\sigma=y\frac{\dd j_0}{j_0}-2\dd y,
 \qquad
 (2\Theta\dd^2\Theta-3(\dd\Theta)^2)|_8\sigma
 =2y\dd^2y-3(\dd y)^2.
\]
If the cusp has width $w$ and local parameter $Q=e^{2\pi i\tau/w}$,
then $j_0=Q^r u(Q)$, $u(0)\ne0$, and
$\dd j_0/j_0=(r+Qu'/u)/w$ has no negative powers.
Both $y$ and its derivatives are holomorphic there. Thus the two
residuals belong to $M_8(\Gamma_0(6))$.

Their coefficients of $q^0,\ldots,q^8$ are zero by the exact product
certificate in Appendix~\ref{app:certificates}.
The index of $\Gamma_0(6)$ in $\mathrm{SL}_2(\Z)$ is
$6(1+1/2)(1+1/3)=12$. The classical valence bound is therefore
$8\cdot12/12=8$. Indeed, for a nonzero holomorphic form of weight $k$ on
$\Gamma_0(6)$, the product of its weight-$k$ slashes over the twelve
cosets is a nonzero holomorphic level-one form of weight $12k$.
Its order at infinity is at least that of the original form, since all
other factors have nonnegative orders. The level-one valence formula
bounds that order by $12k/12=k$. Applying this with $k=8$, neither
residual can have order at least nine unless it is zero. Both residuals
vanish identically.
\end{proof}

\begin{theorem}\label{thm:periods}
The normalized periods satisfy
\begin{equation}\label{eq:periods}
 G(t_+(q))=\Theta_+(q),\qquad F(t_-(q))=\Theta_-(q).
\end{equation}
These are identities of formal series and of their analytic germs at $q=0$.
\end{theorem}
\begin{proof}
Use one row of \eqref{eq:period-rows} and $t=J/\Theta^2$.
All formal substitutions are invertible because $t=q+O(q^2)$.
Set
\[
 e=\frac{\dd t}{t}=\ell-2\frac{\dd\Theta}{\Theta}=\frac W\Theta.
\]
Equation~\eqref{eq:P1} says $e^2=D\Theta^2$, where
$D=1-4at+16ct^2$. The series $e,\Theta,D$ all have constant term one.
Since $\dd=e\theta$, differentiation gives
\[
 2\dd e=(\theta D)\Theta^2+2D\Theta\theta\Theta.
\]
Substitute $\dd\Theta=e\theta\Theta$ and
$\dd^2\Theta=e^2\theta^2\Theta+(\dd e)\theta\Theta$ in
\eqref{eq:P2}. Its left side becomes $\Theta^2I(\Theta)$,
with $I$ as in \eqref{eq:invariant}. Hence $I(\Theta)=0$;
Lemma~\ref{lem:differential} and the fact that $\Theta$ is a formal unit
give $L\Theta=0$.

For each $n\ge1$ the new coefficient of a solution to $LY=0$ is
multiplied by $n^3$, as is explicit in \eqref{eq:central-rec}.
There is a unique solution in $\Q[[t]]$ with constant term one.
The binomially defined series $G,F$ satisfy this equation for the respective
rows and have that constant term. This proves \eqref{eq:periods}.
The eta products and the inverse parameter are analytic near $q=0$, so the
same equalities hold as analytic germs.
\end{proof}

\section{The extended group, divisors and the CM parameter}
\label{sec:geometry}
For a real matrix $M=\smat{a_0}{b_0}{c_0}{d_0}$ with positive determinant,
use the weight-two slash convention
\begin{equation}\label{eq:slash}
 (Y|_2M)(\tau)=\det(M)(c_0\tau+d_0)^{-2}Y(M\tau).
\end{equation}
In particular, if $Y|_2M=\chi(M)Y$ then
\begin{equation}\label{eq:transform}
 Y(M\tau)=\frac{\chi(M)}{\det M}(c_0\tau+d_0)^2Y(\tau).
\end{equation}
Let $\Gamma$ be the projective group generated by $\Gamma_0(6)$ and
\begin{equation}\label{eq:Ws}
 W_3=\smat3163,\qquad W_6=\smat0{-1}60.
\end{equation}

\begin{lemma}\label{lem:character}
The forms $\Theta_\pm$ are modular of weight two on the extended group
$\Gamma$, with characters trivial on their restriction to $\Gamma_0(6)$.
The quotient $\Gamma/\Gamma_0(6)$ has order four projectively.
If $W_2$ denotes the class of $W_6W_3$, the characters are
\begin{equation}\label{eq:characters}
\begin{array}{c|rrrr}
 e&1&2&3&6\\\hline
 \chi_+(W_e)&1&-1&1&-1\\
 \chi_-(W_e)&1&1&-1&-1.
\end{array}
\end{equation}
\end{lemma}
\begin{proof}
The transformation law for eta
\cite[23.18.5--23.18.7]{dlmf}, for
$M=\left(\begin{smallmatrix}a_0&b_0\\c_0&d_0\end{smallmatrix}\right)
\in\mathrm{SL}_2(\Z)$ with $c_0>0$, can be written as
\[
 \eta\!\left(\frac{a_0\tau+b_0}{c_0\tau+d_0}\right)
 =\exp\!\left(\pi i\left(\frac{a_0+d_0}{12c_0}
          -s(d_0,c_0)-\frac14\right)\right)
 (c_0\tau+d_0)^{1/2}\eta(\tau),
\]
with the principal square root and Dedekind sum $s$.
For $\tau'=W_3\tau$ and $L_0=6\tau+3$, the transformations of
$\eta(\tau'),\eta(2\tau'),\eta(3\tau'),\eta(6\tau')$ arise from
\[
 \smat1123,\quad\smat1213,\quad\smat3121,\quad\smat3211
\]
acting respectively on $3\tau,6\tau,\tau,2\tau$.
Their phases, in units of $\pi i$, are
$-1/12,1/12,-1/12,1/12$; the Dedekind sums here are zero.
The phases cancel in each pair, and the linear factors are positive real
multiples of $L_0$, so the principal roots multiply without an extra sign.
Thus
\[
 (\eta_2\eta_3)(\tau')=\frac{L_0}{\sqrt3}(\eta_1\eta_6)(\tau),
 \qquad
 (\eta_1\eta_6)(\tau')=\frac{L_0}{\sqrt3}(\eta_2\eta_3)(\tau).
\]
It follows that $\A(\tau')=L_0^2\B(\tau)/3$ and
$\B(\tau')=L_0^2\A(\tau)/3$. Hence
$\A|_2W_3=\B$ and $\B|_2W_3=\A$.
For $W_6$ the formula $\eta(-1/z)=\sqrt{-iz}\eta(z)$ multiplies each
eta pair by $\sqrt6(-i\tau)$. Both $\A,\B$ acquire the factor
$-6\tau^2$, giving $\A|_2W_6=-\A$ and $\B|_2W_6=-\B$.

The generators normalize the base group. Indeed, for
$\gamma=\smat{a_0}{b_0}{6c_0}{d_0}$,
\begin{align}\label{eq:conjugations}
 W_6\gamma W_6^{-1}&=\smat{d_0}{-c_0}{-6b_0}{a_0},\notag\\
 W_3\gamma W_3^{-1}&=
 \begin{pmatrix}
 3a_0+6c_0-6b_0-2d_0&-a_0-2c_0+3b_0+d_0\\
 6a_0+18c_0-12b_0-6d_0&-2a_0-6c_0+6b_0+3d_0
 \end{pmatrix}.
\end{align}
Moreover
$W_3^2=3\smat5{2}{12}5$, $W_6^2=-6I$ and
$(W_3W_6)^2=-18I$. These relations give at most four classes modulo the
base group and scalars. They are distinct because their determinant square
classes are $1,3,6,2$ in $\Q^\times/(\Q^\times)^2$.
The transformations just proved yield \eqref{eq:characters}.
\end{proof}

Introduce weight-zero modular functions
\begin{equation}\label{eq:h-h3-R}
 h=\frac{\eta_2\eta_6^5}{\eta_1^5\eta_3},\qquad
 h_3=\frac{\eta_3^{12}}{\eta_1^{12}},\qquad R=\A/\B.
\end{equation}
Their characters on $\Gamma_0(6)$ are trivial by the eta criterion.
The cusp orders of $h$ are $(-1,0,0,1)$, and it has no interior poles.
It consequently has one simple pole on the compact curve $X_0(6)$.
The resulting map to $\mathbb P^1$ has degree one, so $h$ is a coordinate
and the function field is $\C(h)$.

\begin{lemma}\label{lem:E}
The following identities hold:
\begin{align}
 h_3(1+9h)&=h(1+8h)^2,\label{eq:E1}\\
 Rh&=(1+8h)(1+9h)=1+17h+72h^2.\label{eq:E2}
\end{align}
\end{lemma}
\begin{proof}
Let $E_1,E_2$ be left side minus right side. Exact product expansion gives
$E_1=O(q^9)$ and $E_2=O(q^{11})$; the coefficient data are in
Appendix~\ref{app:certificates}. Set
\[
 M_1=\eta_1^{17}\eta_3^3,\qquad
 M_2=\eta_1^{17}\eta_3^2\eta_6^7.
\]
Each $M_iE_i$ is a sum of eta monomials with nonnegative exponents, listed
in the appendix. Thus it is holomorphic in $\HH$ and at all cusps.
Since $E_i$ is invariant, the whole sum has the multiplier of $M_i$.
Eta multipliers have order dividing $24$; hence
$(M_1E_1)^{24}\in M_{240}(\Gamma_0(6))$ and
$(M_2E_2)^{24}\in M_{312}(\Gamma_0(6))$ with trivial character.
Their orders at infinity are at least
\[
 24(13/12+9)=242>240,\qquad
 24(65/24+11)=329>312.
\]
The same index-$12$ valence bound forces both forms to vanish.
Since $M_i$ is nonzero in $\HH$, this proves $E_i=0$.
\end{proof}

\begin{proposition}\label{prop:geometry}
Each $t_\pm$ is a Hauptmodul for $\Gamma$.
Each $\Theta_\pm$ has no cusp zeros, and its interior zeros are exactly
the simple zeros above the double poles of $t_\pm$ on $X_0(6)$.
\end{proposition}
\begin{proof}
Equation~\eqref{eq:E2} gives
\begin{align}\label{eq:t-rational}
 t_+&=\frac{R}{(1+R)^2}
 =\frac{h(1+8h)(1+9h)}{(1+6h)^2(1+12h)^2},\notag\\
 t_-&=\frac{R}{(R-1)^2}
 =\frac{h(1+8h)(1+9h)}{(1+16h+72h^2)^2}
 =\frac{t_+}{1-4t_+}.
\end{align}
The numerator and denominator in each expression are coprime.
The quadratic $1+16h+72h^2$ has discriminant $-32\ne0$.
Thus each rational map has degree four and two double poles on $X_0(6)$.
The transformations of $\A,\B$ show that $t_\pm$ are $\Gamma$-invariant.
Since the quotient map has degree four by Lemma~\ref{lem:character},
each $t_\pm$ has degree one on the quotient curve and is a Hauptmodul.

There are no elliptic points for $\Gamma_0(6)$ modulo scalars.
An elliptic matrix of determinant one would have trace $0$ or $\pm1$.
Trace zero implies $a_0^2\equiv-1\pmod3$, impossible; odd trace is
impossible modulo two, since $a_0d_0\equiv1\pmod2$.
Interior local coordinates on $X_0(6)$ are therefore ordinary coordinates.
At cusps $\Theta_\pm$ are nonzero by \eqref{eq:cusp-orders}.
The identity $t_\pm\Theta_\pm^2=J$ and the absence of zeros of $J$ in
$\HH$ imply that above each double pole the corresponding form has a
simple zero, and that it has no other interior zeros.
\end{proof}

\begin{lemma}\label{lem:CM-point}
Let $\delta=\sqrt{-3}$ have positive imaginary part and let
$\beta=\delta/3=i/\sqrt3$. Then
\begin{equation}\label{eq:CM-point}
 h(\beta)=\frac{\sqrt3-1}{24},\qquad
 t_+(\beta)=1/54,\qquad t_-(\beta)=1/50.
\end{equation}
Both $\Theta_+(\beta)$ and $\Theta_-(\beta)$ are nonzero.
\end{lemma}
\begin{proof}
Since $3\beta=-1/\beta$, inversion of eta gives
$h_3(\beta)=(-i\beta)^6=1/27$.
Equation~\eqref{eq:E1} becomes
\[
 27h(1+8h)^2-(1+9h)=(6h+1)(288h^2+24h-1)=0.
\]
Eta products are positive on the positive imaginary axis, so $h(\beta)>0$.
This selects the positive root $(\sqrt3-1)/24$.
The polynomial identity
\[
 54h(1+17h+72h^2)-(1+18h+72h^2)^2
 =-(18h^2-6h-1)(288h^2+24h-1)
\]
and \eqref{eq:t-rational} give $t_+(\beta)=1/54$;
its denominator is nonzero at positive $h$. The last expression in
\eqref{eq:t-rational} yields $t_-(\beta)=1/50$.
Finally $J(\beta)\ne0$ and the two parameters are finite, so neither
form vanishes at $\beta$.
\end{proof}

\section{The Hecke polynomial and finite specialization}
\label{sec:Hecke}
In this section use either pair $(t,\Theta)=(t_+,\Theta_+)$ or
$(t_-,\Theta_-)$, and let $Y=G$ or $F$ respectively.
In particular, $\Theta$ is a weight-two modular form on the extended
group $\Gamma$ with character trivial on $\Gamma_0(6)$, as proved in
Lemma~\ref{lem:character}. This extended-group hypothesis is used in
the Hecke orbit construction \cite[Lemma~3.7]{beukers}.

We spell out that orbit at level six. For $p>3$ set
\[
 \mathcal M(6,p)=
 \left\{\smat{a_0}{b_0}{6c_0}{d_0}:
 a_0,b_0,c_0,d_0\in\Z,\ a_0d_0-6b_0c_0=p\right\}.
\]
The gcd of the first column is $1$ or $p$.
A Bezout identity takes a primitive first column to $(1,0)^T$ by left
multiplication with $\Gamma_0(6)$. If the gcd is $p$, first divide the
column by $p$; its lower entry is still divisible by six.
The $p+1$ distinct left classes therefore have representatives
\begin{equation}\label{eq:Hecke-reps}
 \gamma_\infty=\smat p001,
 \qquad \gamma_k=\smat1k0p\quad(0\le k<p).
\end{equation}
For clarity, when the first-column gcd is $g$, the lower row of the
reducing matrix can be chosen as $(-6c_0/g,a_0/g)$, and a Bezout upper
row makes its determinant one. Its lower left entry is divisible by six
because $p$ is coprime to six. The resulting upper triangular matrix is
$\smat1k0p$ for $g=1$, with $k$ reduced modulo $p$ by a translation;
it is $\smat p001$ for $g=p$, after eliminating its upper right entry.
These representatives are distinct by multiplying the triangular inverses.

The conjugation formulas \eqref{eq:conjugations} preserve
$\mathcal M(6,p)$; they do not require determinant one. Thus $W_3,W_6$
permute its left classes. Put $Q_\gamma=(\Theta|_2\gamma)/\Theta$.
These quotients are independent of the representative because the base
character is trivial. For $w\in\Gamma$ and $\gamma'=w^{-1}\gamma w$,
\[
 Q_\gamma(w\tau)
 =\frac{\Theta|_2\gamma w}{\Theta|_2w}
 =\frac{\Theta|_2w\gamma'}{\chi(w)\Theta}
 =Q_{\gamma'}(\tau).
\]
The character cancels, even though it need not be trivial on $\Gamma$.

\begin{lemma}\label{lem:Hecke-core}
For every prime $p>3$ there is a polynomial
$\Tp(X,t)=\sum_{m=0}^{p+1}C_m(t)X^m\in\Z[X,t]$ such that
\begin{equation}\label{eq:Hecke-product}
 \Tp(X,t(\tau))=
 \left(1-\frac{\Theta(p\tau)}{\Theta(\tau)}X\right)
 \prod_{k=0}^{p-1}
 \left(p^2-\frac{\Theta((\tau+k)/p)}{\Theta(\tau)}X\right).
\end{equation}
Its coefficients satisfy
\begin{equation}\label{eq:Hecke-core}
 \deg C_m\le\lfloor m/2\rfloor,\qquad
 \Tp(X,t)\equiv C_p(t)X^p+C_{p+1}(t)X^{p+1}\pmod{p^2},
 \quad C_{p+1}(t)\equiv1\pmod p.
\end{equation}
\end{lemma}
\begin{proof}
The product in \eqref{eq:Hecke-product} equals
$p^{2p}\prod_{[\gamma]}(1-Q_\gamma X/p)$ using
\eqref{eq:Hecke-reps}. Hence its coefficients are $\Gamma$-invariant.
The numerator of each $Q_\gamma$ is holomorphic in $\HH$.
There are also no poles at cusps: at a cusp represented by $\sigma$,
the quotient is
\[
 \frac{\Theta|_2\gamma\sigma}{\Theta|_2\sigma}.
\]
The denominator has nonzero constant term. Write
$\gamma\sigma=\rho u$, where $\rho\in\mathrm{SL}_2(\Z)$ and $u$
is upper triangular with positive determinant. The numerator is an
expansion of the holomorphic form at a cusp, with the cusp coordinate
rescaled by a positive power. It has no negative exponents.

Thus the only possible poles of a coefficient are the zeros of $\Theta$,
which lie over the pole of $t$. Since $t$ is a coordinate on the quotient,
each coefficient is a polynomial in $t$, initially over $\C$.
A term of the coefficient of $X^m$ has at most $m$ factors $\Theta^{-1}$.
Above $X_0(6)$ its pole order is at most $m$, while the pole order of $t$
is two. A polynomial of degree $d$ in $t$ has pole order $2d$ there.
This proves $\deg C_m\le\lfloor m/2\rfloor$.

For integral coefficients, put $s=q^{1/p}$ and $\zeta=\zeta_p$.
The product expansions lie in $\Z[\zeta][[s]]$.
The symmetry $s\mapsto\zeta s$ removes exponents not divisible by $p$,
and the Galois automorphisms $\zeta\mapsto\zeta^a$ permute the factors.
The surviving coefficients are rational algebraic integers, hence integers
in $\Z[[q]]$. Because $t=q+q^2\Z[[q]]$, successive subtraction of the
lowest terms against powers of $t$ makes all polynomial coefficients
integers. The same triangular argument transfers coefficientwise
congruences from $q$ to $t$.

Let $Q=\Theta(q)/\Theta(q^p)$ and
$R_k=\Theta(q)/\Theta(\zeta^ks)$. The product is exactly
\[
 C_{p+1}(t)(X-Q)\prod_{k=0}^{p-1}(X-p^2R_k).
\]
The symmetric functions of the $R_k$ are in $\Z[[q]]$ by the same
symmetries. Modulo $p^2$, the last product is $X^p$, which proves the
middle congruence in \eqref{eq:Hecke-core}.
The leading coefficient, with positive sign since $p+1$ is even, is
\[
 C_{p+1}(t(q))=
 \frac{\Theta(q^p)\prod_{k=0}^{p-1}\Theta(\zeta^ks)}{\Theta(q)^{p+1}}.
\]
The product in the numerator is integral in $q$.
Modulo $1-\zeta$ it is $\Theta(s)^p\equiv\Theta(q)$ in characteristic
$p$. Since $(1-\zeta)\Z[\zeta]\cap\Z=p\Z$, this gives the congruence
modulo $p$ on its integer coefficients. Finally
$\Theta(q^p)\equiv\Theta(q)^p\pmod p$ proves
$C_{p+1}\equiv1\pmod p$.
\end{proof}

\begin{lemma}\label{lem:endpoint}
For every prime $p>3$,
\begin{equation}\label{eq:endpoint}
 g_p\equiv g_1=6\pmod{p^2},\qquad
 b_p\equiv b_1=4\pmod{p^2}.
\end{equation}
\end{lemma}
\begin{proof}
The exact product
$\binom{2p}p=2\prod_{k=1}^{p-1}(1+p/k)$ and
$\sum_{k=1}^{p-1}k^{-1}\equiv0\pmod p$ give
$\binom{2p}p\equiv2\pmod{p^2}$.
The latter harmonic congruence follows because inversion permutes
the nonzero residues modulo $p$. In $a_p$ all internal terms have the
factor $\binom pk^2$, while the two endpoint terms are
$1$ and $\binom{2p}p$. Thus $a_p\equiv3\pmod{p^2}$.
For $f_p$ all internal terms are divisible by $p^3$, so
$f_p\equiv2\pmod{p^3}$. Multiplication proves \eqref{eq:endpoint}.
\end{proof}

\begin{theorem}[Finite specialization]\label{thm:specialization}
For either pair of period and parameter, the polynomial identity
\begin{equation}\label{eq:finite-specialization}
 C_p(t)+C_{p+1}(t)Y_p(t)\equiv0\pmod{p^2}
\end{equation}
holds in $\Z[t]$. If $t_0\in\Z_p$ and $z$ is an exact root of
$\Tp(X,t_0)$ in a finite extension of $\Q_p$ with valuation zero, then
$Y_p(t_0)\equiv z\pmod{p^2}$ in its ring of integers.
\end{theorem}
\begin{proof}
The series $Q=\Theta(q)/\Theta(q^p)$ is a root of
\eqref{eq:Hecke-product}, and is a formal unit with constant term one.
Dividing its equation from \eqref{eq:Hecke-core} by the formal unit $Q^p$
gives $C_p+C_{p+1}Q\equiv0\pmod{p^2}$.
Write $C_{p+1}=1+pD_p$. Then
\begin{equation}\label{eq:K}
 Q\equiv K(t):=-(1-pD_p(t))C_p(t)\pmod{p^2},\qquad
 \deg K\le(p-1)/2+(p+1)/2=p.
\end{equation}
The formal inverse $q(t)$ is integral. The period identity
\eqref{eq:periods} expresses $Q$ as $Y(t)/Y(t^\sigma)$, where
$t^\sigma=t(q(t)^p)=t^p+O(t^{p+1})$. Hence
\begin{equation}\label{eq:Q-boundary}
 Q=Y_p(t)+(y_p-y_1)t^p+O(t^{p+1}).
\end{equation}
Lemma~\ref{lem:endpoint} makes the coefficient of $t^p$ zero modulo $p^2$.
Both $K$ and $Y_p$ have degree at most $p$, so comparison through degree
$p$ proves their congruence as polynomials. Multiplying by $C_{p+1}$
proves \eqref{eq:finite-specialization}.

We now specialize only these finite polynomials. The exact equation
$\Tp(z,t_0)=0$ and \eqref{eq:Hecke-core} imply
\[
 z^p\bigl(C_p(t_0)+C_{p+1}(t_0)z\bigr)\equiv0\pmod{p^2}.
\]
Divide by the unit $z^p$, subtract \eqref{eq:finite-specialization}
evaluated at $t_0$, and divide by the unit
$C_{p+1}(t_0)\equiv1\pmod p$. The result is the claimed congruence.
\end{proof}

The upper degree $p$ in \eqref{eq:K} is needed: comparison through
$p-1$ alone would not exclude an extra multiple of $t^p$.
Equation~\eqref{eq:endpoint} treats exactly this endpoint.
Every evaluation in Theorem~\ref{thm:specialization} is of a finite
polynomial; it does not require evaluating an infinite series at a
$p$-adic unit. The argument also needs only the modulus $p^2$ in
\eqref{eq:Hecke-core}; the additional hypotheses in the stronger
modulo-$p^3$ results of Beukers are not used.

\section{The split primes and the Franel corollary}\label{sec:split}
Let $p=x^2+3y^2\equiv1\pmod3$ be prime, with $x,y>0$.
Then $\gcd(x,3y)=1$ and $x,y$ have opposite parity.
Keep $\delta=\sqrt{-3}$ and $\beta=\delta/3$.
We first construct exact complex roots of the Hecke polynomials at the
rational parameters in \eqref{eq:CM-point}, then choose their $p$-adic
embeddings.

\subsection{The case \texorpdfstring{$y$}{y} even}
Here $p\equiv1\pmod{12}$ and $\eps=1$.
Choose $a_0,b_0\in\Z$ with $a_0x+3b_0y=1$ and set
\begin{equation}\label{eq:CM-even}
 M=\smat{a_0}{b_0}{-3y}{x}\in\Gamma_0(6),\qquad
 k=b_0x-a_0y.
\end{equation}
The determinant is one and direct rationalization yields
\begin{equation}\label{eq:CM-image-even}
 M\beta=\frac{\beta+k}{p}.
\end{equation}
Indeed the rationalized numerator is
$b_0x-a_0y+(a_0x+3b_0y)\delta/3=k+\beta$.
The character is trivial on $M$, so \eqref{eq:transform} gives
\begin{equation}\label{eq:root-even}
 p^2\frac{\Theta_+(\beta)}{\Theta_+(M\beta)}
 =\frac{p^2}{(x-y\delta)^2}=(x+y\delta)^2.
\end{equation}
This is an exact root of $\Tp(X,1/54)$ by the factor with index
$k\bmod p$ in \eqref{eq:Hecke-product}; period one of $\Theta_+$ justifies
reducing the index. The same equality holds with $\Theta_-$, giving an
exact root of the corresponding polynomial at $1/50$.

\subsection{The case \texorpdfstring{$y$}{y} odd}
Here $x$ is even, $p\equiv7\pmod{12}$ and $\eps=-1$.
Choose $a_0,b_0\in\Z$ with $3a_0y+b_0x=1$ and put
\begin{equation}\label{eq:CM-odd}
 M=\smat{3a_0}{b_0}{-3x}{3y},\qquad \det M=3,
 \qquad k=b_0y-a_0x.
\end{equation}
Its extended-group class is $W_3$, since
\begin{equation}\label{eq:CM-class}
 MW_3^{-1}=
 \smat{3a_0-2b_0}{b_0-a_0}{-3x-6y}{x+3y}\in\Gamma_0(6).
\end{equation}
The lower left entry is divisible by six because $x$ is even.
Again
\begin{equation}\label{eq:CM-image-odd}
 M\beta=(\beta+k)/p.
\end{equation}
Here the rationalized numerator is
$3k+(b_0x+3a_0y)\delta=3k+\delta$, and the denominator norm is $3p$.
Using $\chi_+(M)=1$ and the determinant factor in
\eqref{eq:transform}, we obtain the exact root
\begin{equation}\label{eq:root-odd}
 p^2\frac{\Theta_+(\beta)}{\Theta_+(M\beta)}
 =\frac{3p^2}{(3y-x\delta)^2}
 =\frac{(3y+x\delta)^2}{3}=-(x-y\delta)^2
\end{equation}
of $\Tp(X,1/54)$. For $\Theta_-$ the character on $W_3$ is $-1$,
so the exact root of its polynomial at $1/50$ is instead
$(x-y\delta)^2$.

\subsection{The unit root and completion of the proof}
Since $p\equiv1\pmod3$, $-3$ is a square modulo $p$.
The two roots lift by Hensel's lemma, giving the two embeddings
$\Q(\delta)\hookrightarrow\Q_p$. In either embedding the factors
$x+y\delta$ and $x-y\delta$ are integral and their product is $p$.
They cannot both vanish modulo $p$, for that would force $p\mid2x$ and
$p\mid2y$. Thus their valuations are $0$ and $1$ in some order;
the embeddings interchange that order.

In the even case choose the embedding where $x+y\delta$ is a unit;
in the odd case choose the one where $x-y\delta$ is a unit.
The exhibited root is now a unit before we invoke
Theorem~\ref{thm:specialization}. Its conjugate square is divisible by
$p^2$, and the exact trace identity is
\begin{equation}\label{eq:trace}
 (x+y\delta)^2+(x-y\delta)^2=4x^2-2p.
\end{equation}
Since $1/54\in\Z_p$, Theorem~\ref{thm:specialization} and the roots
\eqref{eq:root-even}, \eqref{eq:root-odd} give
\begin{equation}\label{eq:split-values}
 B_p\equiv
 \begin{cases}
 4x^2-2p,&p\equiv1\pmod{12},\\
 -(4x^2-2p),&p\equiv7\pmod{12}
 \end{cases}\pmod{p^2}.
\end{equation}
Multiplication by $\eps$ proves the split branch of
Theorem~\ref{thm:main}. The inert branch, including $p=5$, was proved
in Section~\ref{sec:invariant}; hence the theorem holds for all $p>3$.

The complex period ratios used here are defined because both forms are
nonzero at $\beta$. The transformations identify those ratios with exact
algebraic numbers. Only those algebraic numbers are subsequently embedded
in $\Q_p$, so the construction requires no $p$-adic evaluation of complex
eta series at the CM point.

The evaluation in \eqref{eq:Franel-value} is the statement of
Sun's 2018 Conjecture~3.5 \cite[Conjecture~3.5]{sunfranel}, distinct from
the 2022 conjecture proved in Theorem~\ref{thm:main}.

\begin{corollary}\label{cor:Franel}
For every prime $p>5$,
\begin{equation}\label{eq:Franel-value}
 T_{50,p}\equiv
 \begin{cases}
 4x^2-2p,&p=x^2+3y^2\equiv1\pmod3,\\
 0,&p\equiv2\pmod3
 \end{cases}\pmod{p^2},
\end{equation}
and
\begin{equation}\label{eq:bridge-corollary}
 B_p\equiv\leg3p T_{50,p}\pmod{p^2}.
\end{equation}
\end{corollary}
\begin{proof}
For split primes the roots for $\Theta_-$ computed above are the unit
square with positive sign in both parity cases. Since $1/50\in\Z_p$,
Theorem~\ref{thm:specialization} and \eqref{eq:trace} give the split value
in \eqref{eq:Franel-value}. Equation~\eqref{eq:inert-result} gives the
inert value. For split primes quadratic reciprocity gives
$\leg3p=(-1)^{(p-1)/2}=\eps$; for inert primes both sums are zero.
This proves \eqref{eq:bridge-corollary} as a consequence of their separate
evaluations.
\end{proof}

\appendix
\section{Exact eta certificates}\label{app:certificates}
This appendix specifies the finite calculations used in
Proposition~\ref{prop:P} and Lemma~\ref{lem:E}.
They certify identities because the proofs place the residuals in spaces
of holomorphic modular forms with explicit valence bounds.

For $E(\tau)=\prod_{d\mid6}\eta(d\tau)^{r_d}$,
the weight is $\tfrac12\sum r_d$. The eta criterion uses
$\sum dr_d\equiv\sum(6/d)r_d\equiv0\pmod{24}$; for the even weights
here, the rational square condition on $\prod d^{r_d}$ gives trivial
character. The cusp order at denominator $c\mid6$ is
\begin{equation}\label{eq:Ligozat}
 \frac6{24\gcd(c^2,6)}
 \sum_{d\mid6}\frac{\gcd(c,d)^2}{d}r_d.
\end{equation}
The data are
\begin{center}\small
\begin{tabular}{c r r r c}
\toprule
Function&$(r_1,r_2,r_3,r_6)$&Weight&$(\sum dr_d,\sum(6/d)r_d)$
 &Cusp orders\\\midrule
$\A$&$(-5,7,7,-5)$&2&$(0,0)$&$(0,1,1,0)$\\
$\B$&$(7,-5,-5,7)$&2&$(24,24)$&$(1,0,0,1)$\\
$J$&$(2,2,2,2)$&4&$(24,24)$&$(1,1,1,1)$\\
$h$&$(-5,1,-1,5)$&0&$(24,-24)$&$(-1,0,0,1)$\\
$h_3$&$(-12,0,12,0)$&0&$(24,-48)$&$(-2,-1,2,1)$\\
$R$&$(-12,12,12,-12)$&0&$(-24,-24)$&$(-1,1,1,-1)$\\
\bottomrule
\end{tabular}
\end{center}
In every row, the valuations at $2,3$ of $\prod d^{r_d}$ are even.

All required expansions use the integer product
\begin{equation}\label{eq:eta-product-algorithm}
 E(q)=q^s\prod_{m\ge1}(1-q^m)^{e_m},\qquad
 s=\frac1{24}\sum dr_d,\qquad e_m=\sum_{d\mid m}r_d.
\end{equation}
For the coefficient of degree at most $N$, only factors with
$m\le N-s$ contribute. Negative powers are expanded with the integer
binomial coefficients for $(1-q^m)^{-r}$.
An equivalent exact recurrence, for
$\prod_{m\ge1}(1-q^m)^{e_m}=\sum_{n\ge0}c_nq^n$, is
\begin{equation}\label{eq:eta-recurrence-algorithm}
 c_0=1,\qquad
 n c_n=-\sum_{j=1}^n\left(\sum_{d\mid j}d e_d\right)c_{n-j}.
\end{equation}

For the period certificate the required coefficients are
\begin{center}\small
\begin{tabular}{r r r r r r}
\toprule
$n$&$[q^n]\Theta_+$&$[q^n]W_+$&$[q^n]\Theta_-$&$[q^n]W_-$&$[q^n]J$\\\midrule
0&1&1&1&1&0\\
1&6&-8&4&-6&1\\
2&6&-40&20&-78&-2\\
3&42&-296&4&-114&-3\\
4&6&-296&52&-654&4\\
5&36&-1008&24&-756&6\\
6&42&-1480&20&-1482&6\\
7&48&-2752&32&-2064&-16\\
8&6&-2344&116&-5262&-8\\\bottomrule
\end{tabular}
\end{center}
Here $W_\pm=\Theta_\pm\dd J/J-2\dd\Theta_\pm$.
Their coefficients can be obtained without division by a series of zero
constant term by first writing $J=q\,j(q)$, $j(0)=1$, so that
$\dd J/J=1+\dd j/j$. In computing through degree eight one expands
$J$ through degree nine. Substitution of the table in
\eqref{eq:P1} and~\eqref{eq:P2}, using
$[q^n]\dd\Theta=n[q^n]\Theta$, gives the zero vector
$(0,0,0,0,0,0,0,0,0)$ for each of the four residuals.

For the geometric certificate the expansions are
\begin{center}\small
\begin{tabular}{r r r r}
\toprule
$n$&$[q^n]h$&$[q^n]h_3$&$[q^n]Rh$\\\midrule
0&0&0&1\\
1&1&1&17\\
2&5&12&157\\
3&19&90&1043\\
4&61&508&5573\\
5&174&2391&25422\\
6&455&9828&102703\\
7&1112&36428&376600\\
8&2573&124188&1275445\\
9&5689&395199&4041161\\
10&12102&1186344&12096534\\\bottomrule
\end{tabular}
\end{center}
These give $[q^0,\ldots,q^8]E_1=0$ and
$[q^0,\ldots,q^{10}]E_2=0$.
After multiplication by $M_1,M_2$ from Lemma~\ref{lem:E}, the monomials
and their coefficients are
\begin{center}\small
\begin{tabular}{c c r}
\toprule
Residual&Exponents $(\eta_1,\eta_2,\eta_3,\eta_6)$&Coefficient\\\midrule
$M_1E_1$&$(5,0,15,0)$&1\\
&$(0,1,14,5)$&9\\
&$(12,1,2,5)$&$-1$\\
&$(7,2,1,10)$&$-16$\\
&$(2,3,0,15)$&$-64$\\\midrule
$M_2E_2$&$(0,13,13,0)$&1\\
&$(17,0,2,7)$&$-1$\\
&$(12,1,1,12)$&$-17$\\
&$(7,2,0,17)$&$-72$\\\bottomrule
\end{tabular}
\end{center}
Nonnegative eta exponents establish holomorphy of these sums at every
cusp. Together with the multiplier argument and bounds $242>240$,
$329>312$, these exact coefficients prove the full identities.

The accompanying software, version~0.1.2, implements the integer products,
the recurrence \eqref{eq:eta-recurrence-algorithm}, and the residual
calculations. The companion repository is
\begin{center}\small
\url{https://github.com/aconsciousfractal/Sun-3-5-Supercongruence}
\end{center}
The repository README and \texttt{REPRODUCE.md} give the verification and
build instructions.
Its additional checks at bounded sets of primes are regression checks;
the general conclusions follow from the arguments in the article.
The evidence inventory is \texttt{EVIDENCE\_SHA256.txt}; its SHA-256 is
\begin{center}\footnotesize\ttfamily
a8bb62720842b534d6d92927c0ac628f33d922eb07d95d642a86d4c8efb63214
\end{center}
This inventory identifies the companion code, tests, configuration and exact
evidence, excluding the manuscript and prose that quote the digest.

\section{Telescoping proofs of the sequence recurrences}
\label{app:recurrences}
For $n\ge2$, regard binomial terms with lower indices outside their
natural range as zero. For the first sequence put
$\tau(n,k)=\binom nk^2\binom{2k}k$ and
\[
 R_a(n,k)=\frac{k^3(3k-4n)}{n^2}.
\]
\begin{samepage}
\allowdisplaybreaks[0]
The ratio
$\tau(n,k+1)/\tau(n,k)=2(2k+1)(n-k)^2/(k+1)^3$ gives the exact identity
\begin{align}\label{eq:telescoping-a}
 &n^2\tau(n,k)-(10n(n-1)+3)\tau(n-1,k)
       +9(n-1)^2\tau(n-2,k)\notag\\
 &\hspace{12mm}=R_a(n,k+1)\tau(n,k+1)-R_a(n,k)\tau(n,k).
\end{align}
\end{samepage}
For example, after dividing by $\tau(n,k)$ the ratios of the two lower
rows are $(n-k)^2/n^2$ and
$(n-k)^2(n-k-1)^2/(n^2(n-1)^2)$, reducing the identity to a polynomial
equality. Summation over $k=0,\ldots,n$ has zero lower boundary because
$R_a(n,0)=0$, and zero upper boundary because $\tau(n,n+1)=0$.
This proves \eqref{eq:recurrence} with $(a,b,c)=(10,3,9)$;
the initial values are $a_0=1,a_1=3$.

For the Franel sequence let $\tau(n,k)=\binom nk^3$ and set
\begin{align*}
 Q_f(n,k)&=4k^3+(6-18n)k^2+(27n^2-15n)k-14n^3+10n^2,\\
 R_f(n,k)&=\frac{k^3Q_f(n,k)}{n^3(n-1)}.
\end{align*}
Now $\tau(n,k+1)/\tau(n,k)=(n-k)^3/(k+1)^3$, and direct polynomial
expansion gives
\begin{align}\label{eq:telescoping-f}
 &n^2\tau(n,k)-(7n(n-1)+2)\tau(n-1,k)
       -8(n-1)^2\tau(n-2,k)\notag\\
 &\hspace{12mm}=R_f(n,k+1)\tau(n,k+1)-R_f(n,k)\tau(n,k).
\end{align}
The lower-row ratios here are $(n-k)^3/n^3$ and
$(n-k)^3(n-k-1)^3/(n^3(n-1)^3)$.
Both boundaries again vanish, so summation proves
\eqref{eq:recurrence} with $(a,b,c)=(7,2,-8)$.
The initial values $f_0=1,f_1=2$ are immediate.
No endpoint certificate has a pole in the summation index: the
denominators displayed above depend only on $n\ge2$.

\section*{AI assistance and responsibility}
AI assistance under the author's direction was used for mathematical
exploration, proof criticism, source comparison, manuscript drafting,
exact computation and repository preparation. OpenAI Codex and GPT-6.1 Sol
assisted with consolidation and separate checks of the proof and companion.
Research supplied through an external agent was also examined; its complete
model provenance is not asserted. The systems had access to related
arguments. Their checks are not independent human peer review, formal
proof verification or a determination of priority. Individual proof steps
are not attributed to particular AI systems. Oleksiy Babanskyy is the human
author and is responsible for the manuscript, source attribution, licence
boundaries and final decision to publish.

\begin{thebibliography}{9}
\bibitem{sun}
Zhi-Hong Sun,
\emph{Supercongruences involving Ap\'ery-like numbers and binomial coefficients},
AIMS Mathematics \textbf{7}(2) (2022), 2729--2781.
\href{https://doi.org/10.3934/math.2022153}{doi:10.3934/math.2022153}.

\bibitem{sunfranel}
Zhi-Hong Sun,
\emph{Congruences for sums involving Franel numbers},
International Journal of Number Theory \textbf{14}(1) (2018), 123--142.
\href{https://maths.hytc.edu.cn/__local/3/45/77/77F7497DDFD531E8688ED13C0D0_262306C8_24EB4.pdf}{Author's full text}.

\bibitem{beukers}
Frits Beukers,
\emph{Supercongruences using modular forms},
arXiv:2403.03301v3 (2025).
\url{https://arxiv.org/abs/2403.03301v3}.

\bibitem{rousewebb}
Jeremy Rouse and John J. Webb,
\emph{On spaces of modular forms spanned by eta-quotients},
arXiv:1311.1460v3 (2014).
\url{https://arxiv.org/abs/1311.1460v3}.

\bibitem{dlmf}
NIST Digital Library of Mathematical Functions,
Section~23.18, \emph{Modular Transformations}, equations 23.18.5--23.18.7.
\url{https://dlmf.nist.gov/23.18}.
\end{thebibliography}
\end{document}
